\begin{textsample}{POS Dim 1 – rr – Score 44.00 – rr/rr\_ef\_61.txt}  \label{ex:f1_pos_217}
5 ORIENTAÇÃO \textbf{METODOLÓGICAS} E PEDAGÓGICAS 5.1 A CONSTITUIÇÃO DO ESPAÇO-TEMPO NAS ESCOLAS DO TERRITÓRIO DE RORAIMA A concepção de espaço e tempo é algo \textbf{inerente} ao ser humano, visto que todas as nossas ações e vivências estão relativamente ligadas a esses fatores, um lugar e um determinado momento.
No percurso de formação e desenvolvimento humano, vários espaços e tempos se constituem, entre estes o espaço e o tempo na educação ; os mesmos são estruturantes do trabalho dentro da escola, pois suas ações ocorrem num espaço ( sala de \textbf{aula}, recreio, quadra, laboratório, biblioteca, pátio, refeitório, sala dos professores, etc. ) e num tempo ( ano letivo, uma semana, uma \textbf{aula} de 50 minutos, uma atividade de 20 minutos, etc. ).
Neste sentido, é salutar perceber a escola como uma instituição que ao longo dos anos foi quebrando determinados paradigmas.
No que concerne à concepção de escola tradicional \textbf{agregou} novos valores e práticas a partir das exigências da sociedade \textbf{contemporânea}.
Cabe \textbf{salientar} que tais mudanças implementadas resultam numa nova \textbf{maneira} de ensinar e, consequentemente, um novo jeito de aprender.
Pensar a escola a partir das necessidades da comunidade escolar, principalmente com o \textbf{foco} no \textbf{aluno} e suas aprendizagens, é também uma \textbf{maneira} de contribuir para a função social que a escola deve exercer na sociedade, enquanto instituição responsável pela formação de cidadãos \textbf{críticos} e, sobretudo, preparados para a vida e os desafios que ela impõe.
Por esse prisma, é cabível refletir sobre a escola e seus espaços ( físico, social, cultural ) e as temporalidades \textbf{inerentes} ( calendário, atividades culturais e outros \textbf{eventos} ), tendo em \textbf{vista} a relevância que cada um ocupa no processo de ensino e aprendizagem.
Para \textbf{tanto}, Assman ( 2007, p. 234 ) acrescenta que “ o bom uso pedagógico do tempo consiste na transformação dos tempos cronológicos, mesmo com a ajuda da compactação desses tempos pela mídia eletrônica, em tempos vivos da experiência do conhecimento ”.
Em outras palavras, as atividades de cunho pedagógico mesmo estando sob controle de um tempo cronológico tem significação muito peculiar, uma vez que são ressignificadas pelas experiências vividas, troca de conhecimento e aprendizado.
Conforme observado por Gallego e Silva ( 2011 ), os tempos e espaços não são neutros, com isso eles passam a ter uma função educativa muito importante.
Desse modo, é \textbf{relevante} refletir sobre os efeitos desses dois aspectos, tempo e espaço, e as marcas deixadas na formação dos \textbf{alunos} e consequentemente no trabalho realizado pelo professor.
Outra discussão bastante \textbf{pertinente}, na revisão dos usos que fazemos dos espaços e tempos escolares, \textbf{diz} respeito ao processo \textbf{avaliativo}.
Gallego e Silva ( 2011 ) \textbf{sugerem} essa revisão a partir de um modelo de avaliação com \textbf{Progressão} Continuada.
Nesse sentido, Gallego e Silva ( 2011, p. 4 ) \textbf{argumentam} que : “ se o \textbf{foco} de nossas preocupações são as aprendizagens, precisamos de um outro tipo de avaliação, mais \textbf{qualitativa}, mais formativa e processual, no sentido de indicar não só os conceitos \textbf{finais}... ” Em outras palavras, a avaliação precisa ser também um indicativo dos avanços e, sobretudo, buscar soluções para os entraves da aprendizagem e alternativas para o desenvolvimento intelectual, social e cultural dos \textbf{alunos}.
Diante desse contexto, inserimos o Território de Roraima, que possui marcas e características próprias \textbf{tanto} nos aspectos físicos e naturais da região, quanto na questão sociocultural.
Pensar a educação no território de Roraima a partir de outras experiências, sem \textbf{perder} de \textbf{vista} suas peculiaridades, é muito enriquecedor e \textbf{sugere} aprimoramento e até mudanças substanciais no modo de conceber a educação e subsidiá -la com práticas inovadoras.
Vale ressaltar que a questão cultural em Roraima é um marco e também uma referência para o Brasil, considerando que Roraima é um dos estados que ainda possui maior número de comunidades indígenas, além das diferentes culturas que aqui se encontram : seja o nordestino que saiu de sua terra, seja o rio-grandense que saiu do outro extremo do Brasil em busca de oportunidades e melhoria de vida e, mais recentemente, a vinda em massa de imigrantes oriundos da Venezuela, também Guyana e Haiti.
As questões mencionadas acima são \textbf{indispensáveis} na \textbf{análise} aqui proposta visando conceituar e refletir sobre o tripé : tempo, espaços escolares e aprendizagem.
Haja \textbf{vista} que as diferentes \textbf{demandas} necessitam de um olhar atento e diferenciado do poder público, que deve adequar todo o sistema de ensino, respeitando a questão cultural, étnica com o \textbf{intuito} maior de atender as \textbf{demandas} com ensino gratuito e de qualidade.
A partir dessas reflexões pode -se inferir que no Território de Roraima, a sua rica diversidade cultural \textbf{influencia} diretamente no calendário escolar, embora saibamos que existam normas e diretrizes que norteiam a organização temporal e espacial das escolas brasileiras, Roraima possui diferenciais quanto sua organização, especificamente nas escolas indígenas e da zona rural.
Como exemplo, temos as manifestações das comunidades indígenas, algo bem específico da cultura local ; as comemorações de datas ligadas ao folclore, bem como as Assembleias dos Tuxauas e reuniões que ocorrem na própria comunidade e outros \textbf{eventos} de cunho cultural local.
As especificidades de cada lugar promovem uma ressignificação do calendário escolar e do espaço também.
A este respeito, Gallego e Silva ( 2011, p. 6 ) afirmam que “ o tempo estrutura a vida social, as instituições e a identidade dos indivíduos, dessa \textbf{maneira}, o tempo consiste num produto de cada sociedade ”.
Destarte, o tempo escolar está diretamente ligado à discussão de cunho administrativo organizacional referente ao calendário escolar ( dias letivos, feriados, férias etc. ), ou seja, um tempo programado para ser executado.
Tal reflexão é \textbf{oportuna} em virtude de questões culturais específicas presente não só no Território de Roraima, mas em outras regiões do país.
A quebra de paradigma provocada por aspectos culturais, por exemplo, \textbf{sugere} uma nova atitude no tratamento com as situações apresentadas e uma delas seria a construção de um Projeto Político Pedagógico efetivado com os membros da comunidade local, promovendo assim, a verdadeira inclusão e valorizando a cultura e as aprendizagens.
No contexto da Educação Especial é necessário abordar que a constituição do espaço e do tempo é também de \textbf{suma} importância.
Os sistemas de ensino devem estar organizados de tal modo, que seus espaços físicos proporcionem acesso aos recursos pedagógicos e de comunicação, eliminando barreiras arquitetônicas e urbanísticas em sua edificação.
O tempo deve respeitar também limites e possibilidades dos \textbf{alunos} inseridos, ao passo que, estes dois aspectos promovam a aprendizagem, valorizando as diferenças e, atendendo assim, as necessidades educacionais de todos os \textbf{alunos}.
Quanto às leis que \textbf{embasam} a educação especial temos a LDB 9394/96, que aponta suportes e equipamentos necessários para atender as necessidades dos \textbf{alunos} com deficiências, transtornos globais do desenvolvimento e altas habilidades/superdotação nas escolas do ensino regular e a Política Nacional da Educação Especial na Perspectiva da Educação Inclusiva ( Artigos 10, 11, 12 e 13 ) também apresentam orientações de como as escolas de ensino regular devem ofertar e se organizar, de modo que todos os \textbf{alunos} tenham suas especificidades atendidas.
Para \textbf{tanto}, ao se pensar no tempo e espaço escolar precisa se levar em consideração a diversidade que envolve as crianças com necessidades especiais.
É importante reafirmar que os espaços da escola podem ser transformados em ambientes pedagógicos e, portanto, podem ser um suporte \textbf{metodológico} para o professor, propiciando a todos os \textbf{alunos} várias possibilidades de aquisição do conhecimento e maior gosto pela escola.
Vale \textbf{salientar} que a reflexão sobre os espaços que, às vezes passam despercebidos aos nossos olhos, mas que podem ser locais de produção e socialização do conhecimento : a quadra de esportes, a biblioteca e o acervo disponível nela, o pátio escolar, o corredor com murais de \textbf{exposições} ou ainda se a escola dispuser de árvores onde pode -se trabalhar roda de leitura ou algum conteúdo voltado para educação ambiental, entre outras atividades.
Por essa \textbf{ótica} a concepção de tempo e espaço no âmbito escolar admite diferentes formas de se conceber essas duas \textbf{categorias} que são atribuições no processo formativo dos \textbf{discentes}.
Para \textbf{tanto}, o Ministério da Educação por meio do Conselho Nacional de Educação corrobora afirmando que : > A organização do percurso formativo, aberto e contextualizado deve ser construída em função das peculiaridades do meio e das características, interesses e necessidades dos estudantes, incluindo não só os componentes \textbf{centrais} obrigatórios previstos na legislação e nas normas educacionais, mas outros também, de modo flexível e variável, conforme cada projeto escolar ( Resolução nº 4, cap. 3º de 13 de julho de 2010 ).
Sobre as garantias e a diversificação do uso do tempo e espaço escolar o mesmo documento destaca : > Concepção e organização do espaço curricular e físico que se imbriquem e alarguem, incluindo espaços, ambientes e equipamentos que não apenas as salas de \textbf{aula} da escola, mas, igualmente, os espaços de outras escolas e os socioculturais e esportivo-recreativos do entorno, da cidade e mesmo da região ( Resolução nº 4, cap. 3º, \textbf{parágrafo} I de 13 de julho de 2010 ).
Outrossim, o conhecimento produzido e compartilhado nos espaços disponíveis da escola pode assumir um caráter interdisciplinar tornando mais enriquecedor a aprendizagem.
Nesse aspecto, a Resolução nº 4, cap. 3º, \textbf{parágrafo} V também assegura : > Organização da \textbf{matriz} curricular entendida como alternativa operacional que embase a gestão do currículo escolar e represente \textbf{subsídio} para a gestão da escola ( na organização do tempo e do espaço curricular, distribuição e controle do tempo dos trabalhos docentes ) passo para uma gestão centrada na \textbf{abordagem} interdisciplinar, organizada por eixos temáticos, mediante interlocução entre os diferentes campos do conhecimento.
Vale ressaltar que a \textbf{abordagem} interdisciplinar já é um esforço visível nas escolas e busca integrar cada vez mais as diferentes áreas do conhecimento como Matemática, Língua Portuguesa, Ciências, Educação Física e outras \textbf{disciplinas} que fazem parte do currículo.
Sendo assim, a consolidação do conhecimento passa a ser uma realidade no cotidiano dos \textbf{alunos}, \textbf{instigando} -os a prosseguirem na busca pelo mesmo e dando maior eficácia ao trabalho dos professores.
Dessa \textbf{maneira}, pode -se \textbf{dizer} que tempo e espaço escolar, enquanto componentes organizacionais das instituições constituem faces da mesma realidade, contudo, a discussão da temática não se esgota nessa \textbf{análise} necessitando de um olhar \textbf{crítico} e de ações concretas voltadas para dentro da escola e que dinamize a transmissão do conhecimento, onde o abstrato e o empírico suscite nos \textbf{alunos} o desejo de conhecer, criticar, pensar e por fim promover as mudanças que a sociedade \textbf{tanto} precisa nos dias \textbf{atuais}.
Superar os antagonismos que rondam a escola e coloca em risco a nobre missão dos educadores, é uma urgência, considerando que a \textbf{teoria} não pode está dissociada da prática do dia a dia.
Tais reflexões convergem para a superação de muitos fatores que ainda são considerados entraves para o êxito escolar : reprovações, evasão escolar, desinteresse dos \textbf{alunos} pela escola e pelo conhecimento e falta de participação efetiva da família no processo ensino-aprendizagem dos filhos.
Por fim, não podermos esquecer -nos da formação continuada dos professores que deve ser oferecida, sendo esta, \textbf{indispensável} para que a administração do tempo e a valorização dos espaços, enquanto ambientes pedagógicos deem sentido ao aprendizado dos \textbf{alunos} e também reconhecimento ao trabalho realizado.

% matched lemmas: abordagem, agregar, aluno, análise, argumentar, atual, aula, avaliativo, categoria, central, contemporâneo, crítico, demanda, discente, disciplina, dizer, embasar, evento, exposição, final, foco, indispensável, inerente, influenciar, instigar, intuito, maneira, matriz, metodológico, oportuno, parágrafo, perder, pertinente, progressão, qualitativo, relevante, salientar, subsídio, sugerir, sumo, tanto, teoria, vista, ótica
\end{textsample}
